\section{Three Numbers, Not One}

\textbf{[STIPULATION]} For every comparison number the series gives three
quantities, never one: the value, the measurement precision of the comparison
quantity, and the determination tolerance of the theory quantity. A percentage
deviation alone says nothing.

The reason is simple and used throughout: a one per cent deviation is meaningless
when the theory quantity is only determined to two significant figures -- and
damning if both sides are fixed to eight figures. Without both tolerances the
percentage is unreadable.

\section{We Compute Nothing}

\textbf{[STIPULATION]} The most important statement about all numbers in this
series: the theory \emph{computes nothing}. It is not a calculation engine
producing predictions to many decimal places. It orders geometric structures and
shows that they fall in the right range.

From this follows how the verification scripts are to be read. When a script
reports agreement ``to $10^{-15}$'', that is \emph{not} a result. It is either
an \textbf{identity} -- a quantity equal by definition or construction ($\alpha = 1$
after charge normalisation, $L_0/\ell_P = \xi$ by definition) -- or it is a
\textbf{range check}: the test of whether we are in the right order of magnitude.
The many decimal places are floating-point arithmetic, not the precision of a
statement.

\textbf{[B]} \emph{Three} types of test, strictly to be distinguished:

An \emph{identity} verifies a definition. It cannot fail as long as the algebra
is correct. Passing it is no achievement; failing it would be a programming error.
Examples: $\alpha = 4\pi/4\pi = 1$ (A130), $L_0 = \xi\,\ell_P$ (A210), the
reverse direction $\xi = 2\sqrt{Gm}$ (A145).

A \emph{range check} compares a structure with measurement and asks only: correct
order of magnitude, correct sign, correct scaling? Whether $E_\text{char}$ is near
$28$ rather than $3$ or $300$, whether $\kappa$ is near $4$ rather than $9$,
whether the CHSH residual is at $10^{-4}$ rather than $10^{-1}$.

A \emph{prediction} is the case where the theory \emph{genuinely} predicts
something measurable: the \textbf{lepton masses}. The structure $(\sqrt2, 2/9)$
fixes $m_\tau$ from $m_e$ and $m_\mu$ -- a sharp, falsifiable prediction,
$m_\tau = 1776{.}968$ MeV against PDG $1776{.}86\pm0{.}12$, i.e.\ $+0{.}90\,\sigma$.
This is not a range check: here a numerical value stands \emph{before} the
measurement, and the measurement decides it.

\textbf{[STIPULATION]} ``We compute nothing'' therefore applies to the
\emph{bookkeeping values} -- $\alpha = 1/137$, the SI value of $G$, the anchor
choice $L_0$ -- but not to the lepton masses. For the bookkeeping values the
question ``is the result exactly correct?'' is wrongly posed; for $m_\tau$ it
is exactly right. The difference: the lepton masses are bound so tightly by a
\emph{ratio} structure that the third mass follows from two -- that is a genuine
prediction, not an absolute value and not a calibration.

\section{The Four Precision Classes of the Series}

\textbf{[B]} Across the whole series four classes appear. Distinguishing them is
the actual work of the tolerance standard.

\textbf{Class 1 -- five figures and more.} The circulant (A110): $Q=2/3$ to six
figures, $\theta(\mu/e)$ to seven. Here a deviation of $10^{-6}$ is a serious
signal, because both sides are fixed that precisely.

\textbf{Class 2 -- per cent level.} The ladder (A100): $\mu/e$ to $0{.}52\,\%$,
$\tau/\mu$ to $1{.}04\,\%$. Here one per cent is \emph{not} a signal against the
structure, because the structure claims nothing more precise at this level -- and
no success that one could report to five figures.

\textbf{Class 3 -- bookkeeping precision.} The SI value of $\alpha$ (A130,
$-0{.}24\,\%$) and of $G$ (A145, $E_\text{char}$ to $1{.}5\,\%$). These
deviations concern the scale conversion, not the structure. They are not a test
(A210); the percentage measures the bridge, not the theory.

\textbf{Class 4 -- undetermined.} The neutrino sector (A150), the quark masses
(A150, scheme-dependent). Here the theory quantity is so little fixed that any
agreement is inconsequential. What cannot be refuted confirms nothing.

\section{The Most Common Errors}

\textbf{[B]} The tolerance standard exists to avoid five errors that all occurred
at least once in the corpus:

\textbf{Confusing classes.} Taking a class-2 deviation (ladder, per cent) against
a class-1 precision (circulant, five figures) -- that is the level error from
A080/A110. Both describe the same masses, but at different levels.

\textbf{Reading bookkeeping as a test.} Counting a class-3 deviation as a
confirmation or refutation. The error A130 names in the $\alpha$ presentation.

\textbf{Counting undetermined as confirmed.} Booking a class-4 quantity as a
success because the theory ``lies in the range'' -- even though the range is so
wide that anything fits.

\textbf{Inventing figures.} Giving a class-2 quantity five significant figures
because the calculation provides them. The floating-point precision of a
calculation is not the determination tolerance of the quantity.

\textbf{Counting dependent as independent.} Presenting a quantity as independent
evidence or as a second figure, even though it \emph{follows arithmetically} from
one already counted. Example from the ladder (A100): in the ratio
$N_1{:}N_2{:}N_3 \approx 1{:}1{.}981{:}2{.}971$, only $N_2=1{.}981$ is measured
-- the third number is forced as $1 + 1{.}981$ and carries no independent evidence.
The standard requires booking such forced figures for what they are -- consequence,
not finding.

\section{The Standard in One Sentence}

\textbf{[STIPULATION]} An agreement counts only as far as the \emph{worse}
determined of the two sides reaches. A deviation weighs only as heavily as the
\emph{better} determined of the two sides supports it.

That is the whole standard. Everything else is its application.

\section{Verification Script}

\begin{itemize}
  \item Classification of the comparison numbers of the series into the four
    classes, with presentation of measurement precision and determination
    tolerance per case; and the countercheck that a class-2 number reported to
    five figures generates spurious precision:
    \texttt{python/A\_Serie\_Skripte/a220\_toleranzmassstab.py}
\end{itemize}

\section{Open Questions}

\textbf{First, the determination tolerance of some theory quantities is itself
not sharp.} For the ladder exponents (A100) it is unclear to how many figures
they are determined at all. Where that is the case, the standard is to be applied
conservatively.

\textbf{Second, the standard does not replace error propagation.} It orders orders
of magnitude; it does not propagate uncertainties through formulas. A full error
propagation the series does not provide.

\section{Supersedes}

This document makes the tolerance rule from A010 and A080 into its own principle
and applies it to all sectors. It subsumes the corresponding parts of Dok.~122
(ratio and absolute) and Dok.~292 (tolerance analysis leptons). Incorporated: R55
(evidential breadth of the ladder coupling: one observation instead of three;
``factor 300'' not as a figure).

\section{Use of AI Tools}

This document was created with the assistance of AI tools. The questions,
stipulations, and all substantive decisions come from the author; the tools
formulate, compute, and create verification scripts.
Responsibility for content and errors rests entirely with the author.
Full wording of the rules in A010.
